\section{Sylow 子群的共轭性}


\textbf{本节用到的记号，先一次交代清楚}：

\begin{itemize}
\item $G$：有限群；$P\in\operatorname{Syl}_p(G)$：一个 Sylow $p$-子群；$Q\le G$：任意 $p$-子群。
\item $\operatorname{Syl}_p(G)$：全体 Sylow $p$-子群；$G$ 通过共轭作用在它上面，$g\cdot P=gPg^{-1}$。
\item $N_G(P)=\{g\in G:gPg^{-1}=P\}$：$P$ 的正规化子，它正是上述作用中 $P$ 的稳定子。
\item $C_G(P)$：$P$ 的中心化子。注意：共轭作用的稳定子是 $N_G(P)$ 而\textbf{不是} $C_G(P)$；只有 $P$ 是 Abel 群时二者才相同。
\item $|G|=p^am,\ p\nmid m$：阶的分解；$[G:N_G(P)]$：指数。
\item $n_p=|\operatorname{Syl}_p(G)|$：Sylow $p$-子群的个数。
\item $(c_1\ c_2\ \cdots\ c_k)$：轮换记号；$S_n$：对称群；$\trianglelefteq$：正规子群；$\le$：子群。
\end{itemize}


\subsection{第二定理的内容}


设 $|G|=p^a m$，其中 $p\nmid m$，并令 $P\in\operatorname{Syl}_p(G)$。Sylow 第二定理包含两句话：

1. 任意 $p$-子群 $Q\le G$ 都包含在某个 $P$ 的共轭子群中，即存在 $g\in G$ 使

\[
   Q\le gPg^{-1}.
\]

2. 任意两个 Sylow $p$-子群彼此共轭。特别地，它们同构，但通常不是同一个子群。

第一句话更强：它告诉我们任意 $p$-子群都可以“放大”到某个 Sylow 子群里；第二句话是取 $|Q|=p^a$ 得到的。

\subsection{用左陪集证明包含结论}


令 $Q\le G$ 是任意 $p$-子群。考虑左陪集集合

\[
\Omega=G/P=\{gP:g\in G\}.
\]

这里 $G/P$ 只是左陪集的集合，\textbf{暂时不能当作商群}，因为 $P$ 未知是否正规。陪集个数为

\[
|\Omega|=[G:P]=m,
\]

不被 $p$ 整除。$Q$ 通过左乘作用在 $\Omega$ 上：

\[
q\cdot(gP)=(qg)P.
\]

这个确实是群作用，因为左乘保持陪集的形式。由轨道--稳定子定理，每个 $Q$-轨道的大小都是 $|Q|$ 的因子，因而是 $p$ 的幂。把 $\Omega$ 分解为不相交轨道后，若所有轨道大小都被 $p$ 整除，则 $|\Omega|$ 也被 $p$ 整除，矛盾。因此存在一个大小不被 $p$ 整除的轨道；由于它的大小是 $p$ 的幂，只能是 $1$。

所以有某个 $gP\in\Omega$ 被 $Q$ 的每个元素固定：

\[
qgP=gP\qquad(\forall q\in Q).
\]

左乘 $g^{-1}$，得到

\[
g^{-1}qgP=P.
\]

一个左乘 $hP$ 等于 $P$ 当且仅当 $h\in P$，故 $g^{-1}qg\in P$。对所有 $q\in Q$ 成立，于是

\[
g^{-1}Qg\le P,
\qquad\text{即}\qquad
Q\le gPg^{-1}.
\]

这就证明了：每个 $p$-子群都包含于某个 Sylow $p$-子群的共轭中。

\subsubsection{等价的“交”形式（[DF] 的 Lemma 19）}


上面的结论有一个看起来不同、其实等价的写法，它在 [DF] 里是证明 Sylow 定理的起点：

\textbf{引理。} 设 $P\in\operatorname{Syl}_p(G)$，$Q\le G$ 是任意 $p$-子群，则

\[
Q\cap N_G(P)=Q\cap P.
\]

\textbf{证明。} 因为 $P\le N_G(P)$，所以 $P\cap Q\le N_G(P)\cap Q$，只需证反向包含。取 $x\in Q\cap N_G(P)$。由 $x\in N_G(P)$ 知 $xPx^{-1}=P$，于是 $xPx^{-1}\le P$；而 $x\in Q$，所以 $xPx^{-1}\le Q$。两边合起来

\[
xPx^{-1}\le P\cap Q.
\]

取阶得 $|P|\le|P\cap Q|$，结合 $P\cap Q\le P$ 得到 $P\cap Q=P$，即 $P\le Q$。又 $x\in Q$ 且 $P\le Q$，所以 $x\in P$。这就证明了 $Q\cap N_G(P)\le Q\cap P$。$\square$

两种写法的关系是直接的：$Q\le gPg^{-1}$ 等价于 $g^{-1}Qg\le P$，也就是“某个共轭的 $Q$ 被 $P$ 吸收”；而引理说的是“$Q$ 中那些把 $P$ 保持为集合的元素，本来就已经在 $P$ 里”。读 [DF] 时如果对 Lemma 19 的作用感到突然，可以回到上面陪集作用的版本：两者证的是同一件事。

\subsection{两个 Sylow 子群必共轭}


若 $Q$ 本身也是 Sylow $p$-子群，则 $|Q|=|P|=p^a$。上面的包含关系给出

\[
Q\le gPg^{-1},\qquad |Q|=|gPg^{-1}|.
\]

有限群中阶相同的包含子群只能相等，所以 $Q=gPg^{-1}$。因此

\[
\operatorname{Syl}_p(G)=\{gPg^{-1}:g\in G\},
\]

是 $G$ 在共轭作用下的一个轨道。

\subsection{正规化子：不要和中心化子混淆}


对子群 $P\le G$，定义正规化子

\[
N_G(P)=\{g\in G:gPg^{-1}=P\}.
\]

它是 $G$ 的子群，并且 $P\le N_G(P)$。由定义，$P$ 在 $N_G(P)$ 中正规：

\[
P\trianglelefteq N_G(P).
\]

这里有两个相似但不同的对象：

\begin{itemize}
\item 元素 $x$ 的中心化子 $C_G(x)=\{g:gx=xg\}$，要求每个元素逐点固定 $x$；
\item 子群 $P$ 的正规化子 $N_G(P)=\{g:gPg^{-1}=P\}$，只要求整体保持子群 $P$ 不变。
\end{itemize}


在 $G$ 共轭作用于 $\operatorname{Syl}_p(G)$ 时，$P$ 的稳定子正是 $N_G(P)$，不是 $C_G(P)$。课堂转写中“除以中心化子”的地方应改为“除以正规化子”。

\subsection{唯一性与正规性}


以下三个条件等价：

\[
n_p=1
\quad\Longleftrightarrow\quad
P\trianglelefteq G
\quad\Longleftrightarrow\quad
gPg^{-1}=P\ (\forall g\in G).
\]

证明很短。若 $n_p=1$，那么 $gPg^{-1}$ 仍是 Sylow $p$-子群，只能等于唯一的 $P$，所以 $P$ 正规。反过来若 $P$ 正规，则所有共轭 $gPg^{-1}$ 都等于 $P$；而第二定理说所有 Sylow 子群都来自这些共轭，因此只有一个。

\textbf{“唯一”指的是什么？} 这里的唯一是集合意义上的“只有一个子群”。一般情况下不同 Sylow 子群虽不相等，却都共轭；说“唯一到共轭”则是说同一个素数对应的 Sylow 子群只有一个共轭类，这个结论在第二定理中总是成立。

\subsection{例：$S_4$ 的 Sylow $2$-子群}


\[
|S_4|=24=2^3\cdot3,
\]

所以 Sylow $2$-子群阶为 $8$。一个具体构造是把两个二元块 $\{1,2\}$ 与 $\{3,4\}$ 互换：块内的交换 $(1\,2)$、$(3\,4)$，以及交换两块的置换 $(1\,3)(2\,4)$ 共同生成一个阶为 $8$ 的子群。任意另一个阶为 $8$ 的子群都与它共轭，但未必等于这个具体子群。

\subsection{参考对照}


\begin{itemize}
\item [DF] §4.5，Theorem 18(2) 及其陪集/轨道证明。
\item [DS] 该书没有正规化子的讨论；正规化子的定义与 "$P\trianglelefteq N_G(P)$" 来自 [DF] §4.4 及本节给出的定义。
\end{itemize}
